\section{Bunch-crossing assignment: checks}
\label{sec:crossing}

\subsection{Why the crossing matters}
In streaming readout one event (one GL1 trigger) contains tracks from many collisions in different
bunch crossings (\vtxPerEv{} vertices per event, Figure~\ref{fig:event}). The silicon detectors
fix the crossing of a track (INTT timing); the TPC measures only the drift time, so the $z$
position of the TPC part of a track depends on the crossing assumed for it. Combining two tracks
from different crossings, or attaching a candidate to a PV from another crossing, produces
combinatorial background and wrong pointing/decay-length values. The default reconstruction
therefore requires both daughters and the PV to have the same crossing
(Section~\ref{sec:reco}).

The vertex-crossing distribution (Figure~\ref{fig:event}, left) is flat at
$\approx$\vtxPerOtherX{} vertices per crossing, with a spike of \vtxAtXzero{} vertices at $X=0$;
\evVtxXzeroFrac\% of the events have a vertex at $X=0$ (\evVtxXzeroMultiFrac\% more than one).
We interpret the spike as the collision that fired the trigger; this interpretation was not
checked against the trigger information. Because all TPC-only tracks have $X=0$, a candidate
with a TPC-only daughter can only be formed with $X=0$ silicon tracks and an $X=0$ PV.

\subsection{Consistency of the default sample}
In the default \Ks{} and \Lam{} ntuples every candidate satisfies
\fcmd{track_1_bunch_crossing} = \fcmd{track_2_bunch_crossing} = \fcmd{primary_vertex_ID}
(violations: \nViolXKs{} and \nViolXLambda; check \checkXKs/\checkXLambda,
\fpath{scripts/make_numbers.C}). \nTTKs{} \Ks{} and \nTTLambda{} \Lam{} candidates have two
TPC-only daughters (before QA cuts).

\subsection{Duplicates across events}
\label{sec:dup}
Consecutive streaming events overlap in time, so the same collision can be reconstructed twice,
with crossings that differ by the BCO difference of the two events. Such candidates have the same
\fcmd{Collision_BCO} and identical kinematics. Using the key (\fcmd{Collision_BCO}, mass rounded to
10\,keV, $p_x$ of \fcmd{track_1} rounded to 0.1\,MeV) we find \dupFracKs\% (\Ks) and
\dupFracLambda\% (\Lam) duplicates before QA cuts, \dupQAKs{} and \dupQALambda{} after QA cuts. They
are removed in all results of this note.

\choice{Duplicates are kept in the ntuple and removed offline with the \fcmd{Collision_BCO} key: the
removal needs the neighbouring events, which a per-event module cannot see across file
boundaries.}

\subsection{Reconstruction without any crossing requirement}
To test whether the crossing requirement loses real signal, the \fcmd{_noXmatch} instances run with
the same cuts but with \fcmd{requireBunchCrossingMatch(false)} and
\fcmd{requireTrackVertexBunchCrossingMatch(false)}: any two tracks may be combined and the PV is
chosen among all vertices by the smallest $\chi^2$-deviation. (The implicit requirement that each
daughter has a vertex in its own crossing, Section~\ref{sec:trackcuts}, still applies.) This gives
\candPerEvNoXKs{} \Ks{} and \candPerEvNoXLambda{} \Lam{} candidates per event. Each candidate is
put in one of six classes:
\begin{itemize}\setlength\itemsep{0pt}
\item[A] both daughters with silicon, $X_1=X_2=X_{\rm PV}$ (what the default keeps);
\item[B] both with silicon, $X_1=X_2\neq X_{\rm PV}$ (consistent pair, PV of another crossing);
\item[C] both with silicon, $X_1\neq X_2$;
\item[D] one TPC-only daughter ($X=0$) and one silicon daughter with $X\neq0$;
\item[E] one TPC-only daughter and one silicon daughter with $X=0$;
\item[F] two TPC-only daughters.
\end{itemize}
The signal in each class is obtained with the QA cuts and a fit with mean and width fixed to the
default-sample result (Table~\ref{tab:xclass}, Figures~\ref{fig:xclassKs}--\ref{fig:xclassLambda}).

\begin{table}[H]
\centering\small
\caption{Signal per crossing class in the no-crossing-requirement samples (QA cuts; Gaussian
with fixed mean/width + free pol2; statistical errors). Duplicates are not removed here.}
\label{tab:xclass}
\begin{tabular}{llrrr}
\toprule
\multicolumn{5}{c}{\Ks} \\
class & definition & candidates & signal & signal/error \\
\midrule
\input{tables/tab_xclass_Ks.tex}
\midrule
\multicolumn{5}{c}{\Lam+\aLam} \\
class & definition & candidates & signal & signal/error \\
\midrule
\input{tables/tab_xclass_Lambda.tex}
\bottomrule
\end{tabular}
\end{table}

\paragraph{Observations.}
\begin{itemize}\setlength\itemsep{1pt}
\item \textbf{Class C} (daughters in different crossings) contains no significant signal:
  \Ks{} $\xclCKs\pm\exclCKs$, \Lam{} $\xclCLambda\pm\exclCLambda$ ($2.6\sigma$), on a background
  of $>3\times10^4$ candidates. Splitting class C, the \Lam{} excess sits entirely at
  $|X_1-X_2|>5$ ($\xclClargeLambda\pm\exclClargeLambda$; $|X_1-X_2|\leq5$:
  $\xclCsmallLambda\pm\exclCsmallLambda$), where the background is largest; for \Ks{} both parts
  are compatible with zero ($\xclCsmallKs\pm\exclCsmallKs$, $\xclClargeKs\pm\exclClargeKs$).
\item A fit-free cross-check: for silicon--silicon pairs, peak window ($|m-\mu|<2\sigma$) minus
  equal-width sidebands (5--7$\sigma$) gives \Ks{} $\dXzeroAllKs\pm\edXzeroAllKs$ with
  $X_1=X_2$ and $\dXnonzeroAllKs\pm\edXnonzeroAllKs$ with $X_1\neq X_2$; for \Lam{}
  $\dXzeroAllLambda\pm\edXzeroAllLambda$ and $\dXnonzeroAllLambda\pm\edXnonzeroAllLambda$
  (Figure~\ref{fig:dX}, which shows $|X_1-X_2|\leq20$).
\item \textbf{Class B} holds $\xclBKs\pm\exclBKs$ \Ks: same-crossing pairs for which the
  smallest-$\chi^2$ PV is from another crossing. Without the crossing requirement the PV choice
  is therefore wrong for about 10\% of the real \Ks{} (B/(A+B)), and the pointing variables are computed
  w.r.t.\ the wrong vertex. In the default sample these \Ks{} are attached to the PV of their own
  crossing: the default yields more \Ks{} with two silicon daughters (sum of the II, IM and MM rows of
  Table~\ref{tab:pairtype}) than class A ($\xclAKs\pm\exclAKs$), presumably because pointing cuts evaluated
  w.r.t.\ a wrong PV reject part of the signal.
\item \textbf{Class D} (TPC-only daughter + silicon daughter with $X\neq0$) has no signal
  (\Ks{} $\xclDKs\pm\exclDKs$, \Lam{} $\xclDLambda\pm\exclDLambda$): with the wrong crossing the
  TPC-only track does not form a good vertex. Such decays are lost both with and without the
  crossing requirement.
\item Classes E and F (crossing 0) carry small signals, consistent with the TT/IT entries of
  Table~\ref{tab:pairtype}.
\end{itemize}

\choice{Keep the default: both daughters and the PV in the same crossing. Relaxing it adds class C,
which after QA cuts has about 16 times more \Ks{} candidates than class A (Table~\ref{tab:xclass})
but no significant signal, and it makes the PV choice wrong for about 10\% of the \Ks{} (class B).}

\caveat{The \Lam{} class-C excess of $\xclCLambda\pm\exclCLambda$ is not significant but not zero
either. If real, it would be \Lam{} whose daughters were given different crossings (e.g.\ a
mis-assigned crossing of a low-momentum proton). This should be rechecked with $\sim$10 times more
events and with the crossing difference resolved per daughter type before it is ignored.}

\begin{figure}[H]
\centering
\includegraphics[width=0.48\textwidth]{fig_dX_Ks}\hfill
\includegraphics[width=0.48\textwidth]{fig_dX_Lambda}
\caption{Crossing difference $X_1-X_2$ for silicon--silicon pairs in the no-crossing-requirement
samples with QA cuts. Black: peak window $|m-\mu|<2\sigma$; red: sidebands 5--7$\sigma$ (same total
width); blue points: difference. Left \Ks, right \Lam. The subtracted signal sits at $X_1=X_2$.}
\label{fig:dX}
\end{figure}

\begin{figure}[H]
\centering
\includegraphics[width=0.95\textwidth]{fig_xclass_Ks}
\caption{$\pi^+\pi^-$ mass per crossing class A--F (no crossing requirement, QA cuts); blue: fit with
mean and width fixed to the default-sample peak. Numbers in Table~\ref{tab:xclass}.}
\label{fig:xclassKs}
\end{figure}
\begin{figure}[H]
\centering
\includegraphics[width=0.95\textwidth]{fig_xclass_Lambda}
\caption{As Figure~\ref{fig:xclassKs} for $p\pi$.}
\label{fig:xclassLambda}
\end{figure}
